\documentclass[11pt,a4paper]{article}
\usepackage[english]{babel}
\usepackage[T1]{fontenc}
\usepackage{lmodern}
\usepackage[a4paper,margin=1in]{geometry}
\usepackage{color} % Define colors for hyperlinks
\usepackage{graphicx}
\usepackage{float}
\usepackage{booktabs}
\usepackage{array}
\usepackage{caption}
\usepackage{subcaption}
\usepackage{tabularx}
\usepackage{amsmath}
\usepackage{hyperref}
\usepackage{enumitem}
\usepackage{titlesec} % Include the titlesec package
\usepackage[nottoc]{tocbibind} % Manage Table of Contents entries
\usepackage[authoryear]{natbib}
\usepackage{fancyhdr}
\usepackage{titlesec}
\usepackage{enumitem}
\usepackage{xcolor}    
\usepackage{pdflscape}  % For landscape pages
\usepackage{xfrac}
\usepackage{tikz}         % For drawing with TikZ
\usetikzlibrary{arrows.meta, positioning, shapes.geometric}
\usepackage{amsmath}
\usepackage{amssymb}
\usepackage{bm}
 
\usepackage{pdflscape}  % For landscape pages


\definecolor{darkblue}{rgb}{0,0,0.8}
\hypersetup{
	colorlinks=true, 
	citecolor=darkblue,
	linkcolor=darkblue, 
	urlcolor=darkblue
}

% ---------- Optional: cleaner paragraphs ----------
\setlength{\parskip}{0.55em}
\setlength{\parindent}{0pt}

\begin{document}
	

\tableofcontents

\pagebreak

\section{Introduction}

\section{Historical Trace of Capacity Utilization Measurement throughout Fordist to Post-Fordist era} 

This section reconstructs a historical trace of capacity utilization measurement not as a narrow sequence of statistical refinements, but as a history of state-crafting of capitalist development adopting operational routines of macroeconomic measurement within and across different institutional configurations looking at the case of the United States. Under a regulationist approach, here is argued that during the Fordist era, macroeconomic measurement of capacity utilization as a practice learned to make slack legible, actionable, and contestable. Its central premise is that operational routines are not neutral tools. Rather, they are institutionalized forms of ideological embeddedness shaping the institutional capacities to implement macroeconomic policy as \textit{spatio-temporal fixes} by means of strategic selectivities on what counts as the relevant signal, which institutions are authorized to act on it, and over what temporal horizon they are expected to intervene. For that same reason, however, they can also become sources of dysfunction when the contradictions they once regularized deepen or shift. What matters, then, is not measurement practices themselves, but the political economy through which they are constructed, interpreted, stabilized, repurposed, and contested across changing accumulation regimes. In \citet{Jessop2003}'s terms, these routines belong to spatio-temporal fixes that secure only provisional coherence and begin to unravel when the ideological commitments in which they are embedded are no longer capable of providing institutional coherence by means of displacing sources of stagnation and crisis junctures to other geographies or future temporal horizons \citep{Jessop2006}. 

This perspective an historical anchor preceding the rise of American Keynesianism, departing by the earliest large-scale attempt to measure the slack of productive capacities in the United States. Interestingly, these efforts did not first emerge from a mature macroeconomic policy apparatus, but from the research program at the Brookings Institution, a long-standing policy think-tank within the landscape of private institutions shaping the \textit{"policy"} debate in Washington, DC. In the aftermath of the Great Depression during the mid-1930s, Brookings' economists expanded empirical economic inquiry studying the production, consumption, capital formation, and income, at national scale, presenting the first empirical efforts to estimate the nation’s productive ceiling sector-by-sector. By doing so, the debate and engagement with macroeconomic policy gain paramount self-reflexive capacities developing critical diagnosis identifying the very source of the great depression,  \emph{idle capacities}, into a measurable ceiling problem \citep{Nourse1934}. That shift pioneered later efforts to establish a policy  dashboard of a broader set of operational routines. Long before American Keynesian became mainstream, \textit{"America’s Capacity to Produce"} -- and its companions on issue of consumption among others, by the Brookings
Institution made visible what counted as idle capacities in the first place. Capacity utilization should therefore not be approached as the progressive discovery of a timeless measurement artifact. It was assembled historically through institutional inquiry, statistical construction, and policy need. From the outset, the measurement issue was inseparable from the problem of governability.

The section follows that object across two historically distinct regime settings. Under Fordism, survey-based utilization measures were stabilized within a broader settlement in which American Keynesianism, business-information routines, and official statistical authority made slack governable for demand management. The crisis of the 1970s did not abolish those routines, but recoded them in a process that unfolded 
as a gradual institutional change \citep{MahoneyThelen2009}. Measures that had helped organize expansionary intervention became increasingly read as warnings of inflationary pressures, while not here demonstrated, an interesting feature of this debate in historical perspective, is the participation of some of the main characters of neoliberal policies in the process of recodification from measuring capacity utilization as a slack index to an inflation sentinel as for example Alan Greenspan.  

The measurement of capacity utilization during the post-fordist era can be interpreted as a discontinuity of the operational routines conducting monetary and fiscal policy, which became clear only ex-post it happened. Specifically in the aftermath of the so-called \textit{"Volcker shock"}, amidst  the rise of econometrics in the field of economics. Particularly, with the rising traction of time series econometrics applications to macroeconomics  within the discipline. \citet{Jessop2002} frames this shift as a reorganization of the forms through which accumulation was regularized: where Atlantic Fordism gave pride of place to wage and money forms institutionalized in the Keynesian National Welfare State (KNWS), the  post-Fordist process of capitalist re-structuring which included the dismantling of the KNWS amidst new institutional forms of competition and correspondingly new ideological commitments the imperative functions of the state. \citet{Vidal2013} places that reorganization in its historical setting by linking the crisis of Fordism to declining profitability, intensified competition, overaccumulation, and the rise of finance. In that setting, is the recodification of slack into a portable statistical device re-conceptualized from capacity utilization to output-gap from  potential-output aligned with an anti-inflationary form of governance.

This shift also clarifies the stakes for the debates around business cycles and capacity utilization for researchers in political economy. Under a theoretical commitment of demand-led growth and distributive conflict, once money is understood as endogenous, the rate of interest can no longer be read as a market-clearing price but as a variable of distributive conflict where low-inflation is favourable for the banking-industry \citep{RochonSetterfield2008,LofaroMatamorosRochon2025}. The institutionalization of interest-rates and inflation targeting as a policy variable set by central banks, then can be understood as an historically specific objective: inflation control while narrowing the space for contested forms of macroeconomic management. This is an essential axis on the rise of finance capital as the hegemonic force within the historical block of the transnational capitalist class \citep{Robinson2000}. 

The stakes are not merely alternative traditions' preference for different
measures. Capacity utilization measurements carry distinctive ideological
commitments institutionalized in operational routines and bureaucratic
practices, presented as uncontestable by means of their depoliticization
\citep{Best2025}. Their construction, sedimentation, and recodification
across accumulation regimes is the object of the historical trace that
follows.

\subsection{From Survey Measures to an Inflation Sentinel: Capacity Utilization Measurements during the Fordist Era}
\label{subsec:21_fordist_capacity_utilization}

The Fordist history of capacity utilization measurement begins with a problem of governability rather than with a settled variable. In the United States, the Great Depression made idle resources visible as a political and administrative problem, but not yet as a stabilized object of macroeconomic management. What had to be built first was a way of rendering slack statistically legible. Brookings-era research was decisive because it turned unused productive capacity from an impressionistic sign of downturn into a measurable ceiling problem \citep{Nourse1934}. The issue was no longer simply whether factories stood idle, but how far actual output fell short of a plausible productive limit and under what conventions that limit should be defined. Capacity utilization thus entered macroeconomic discourse as an operational problem of visibility: institutions had to decide what counted as ``capacity,'' how it could be approximated, and how that approximation could orient intervention. Later assessments made clear that this early construction of the object rested on unresolved tensions between engineering and economic notions of capacity \citep{Burns1954,KleinDavid1958}. The early history of utilization measurement was therefore already a history of governance by numbers rather than a neutral exercise in statistical description.

American Keynesianism gave that problem a clearer institutional and theoretical function. Hansen’s interwar and wartime writings help explain why slack could no longer remain a residual category once stabilization became a policy objective \citep{Hansen1936,Hansen1941}. Hicks and Samuelson then carried that problem into a more formal and pedagogical macroeconomic grammar, helping turn slack into a governable macroeconomic object rather than a residual business-cycle symptom \citep{Hicks1937,Samuelson1947,Samuelson1948}. Once the state was expected to manage demand, sustain employment, and monitor bottlenecks, it needed indicators that translated underutilized productive resources into signals for fiscal and monetary action. Capacity utilization became necessary not because Keynesianism discovered a timeless economic magnitude, but because a full-employment policy regime required an index through which effective-demand deficiencies could be rendered legible to authorities charged with stabilization. That shift mattered because it moved utilization measurement from descriptive inquiry into the core of a policy dashboard. The state did not govern an abstract ``economy'' in general. It governed through a portfolio of indicators that identified idle resources, pressures, and thresholds for intervention. In that setting, utilization was never merely technical. It was bound to a historically specific settlement in which demand management, employment commitments, and macroeconomic expertise reinforced one another \citep{Barber1987}.

That theoretical consolidation did not by itself settle the routine. McGraw--Hill did much of that work. Its plant-and-equipment surveys regularized managerial reports of actual and ``preferred'' operating rates and, in doing so, converted heterogeneous notions of productive capacity into a repeated business-information device \citep{Butler1958,Phillips1963}. What mattered here was less the elimination of definitional ambiguity than the creation of an administratively reproducible benchmark. Under Fordism, reproducibility itself became a source of authority. Repeated survey collection generated comparability across time and sectors, while the broader information infrastructure of monopoly capital helped stabilize the measure’s legitimacy inside business forecasting and policy discussion \citep{BaranSweezy1988,Munroe2007}. What emerged, then, was not simply a better estimate of productive ceilings, but a routinized operational device through which firms, analysts, and policymakers could speak a common language about slack. At the same time, the benchmark’s apparent neutrality rested on a dense set of conventions: assumptions about maintenance, overtime, sustainable operation, and the reporting practices of firms. Survey settlement therefore did not abolish mediation. It normalized it.


The Federal Reserve then transformed that settled routine into official macroeconomic authority. Once capacity and utilization measures were incorporated into the Industrial Production program and disseminated through the G.17 release, the survey benchmark ceased to be merely one possible way of gauging slack. It became the official routine through which the state rendered productive underuse visible \citep{Board2017}. That move altered the structure of the field. Institutional embedding did not simply add prestige to an existing measure; it raised the switching costs of alternatives and standardized what counted as admissible evidence. Physical-intensity measures such as Electric Motors Utilization and the Average Workweek of Capital remained conceptually important and in some respects tracked operational realities more directly, but they failed to secure comparable institutional adoption \citep{Foss1963,TaubmanGottschalk1971}. The point is not that every rival was empirically defeated. It is that one routine proved easier to reproduce administratively, easier to integrate into policy practice, and easier to stabilize inside the state’s statistical machinery. In this way, capacity utilization became an official governance variable of Fordist macroeconomic management \citep{Jessop2003}.

The crisis of the 1970s did not abolish that routine. It recoded it. Under stagflationary pressure, capacity utilization was increasingly detached from its earlier role as a signal of effective-demand slack and reinterpreted as a warning device for inflationary strain. The benchmark survived, but the admissible meaning of the benchmark narrowed. What had once helped justify expansionary intervention now mattered insofar as it signaled when output was approaching inflationary constraint. This shift appears clearly in the Brookings debates, where capacity measures were increasingly judged by whether they could register inflationary pressure credibly \citep{KleinEtAl1973}. This was not a minor technical refinement. It marked a crisis-driven transformation in the policy dashboard itself. As the Bretton Woods order unraveled, the Fordist compromise weakened, and profitability pressures intensified, the utilization index was repurposed within a broader reorientation of macroeconomic governance. The same operational routine that had once helped organize full-employment management was now made compatible with emerging anti-inflation discipline. In regulationist terms, the benchmark persisted across a transition between modes of regularization, but its place in the hierarchy of policy priorities changed decisively \citep{Aglietta2000,Jessop2002}. That is why the recoding of capacity utilization should be read as a critical juncture in the fall of American Keynesianism rather than as a simple update in measurement technique.

What survived this transition was not Keynesian slack governance itself, but the institutional durability of the routine. Survey-based utilization retained authority because it had already been stabilized inside the policy apparatus, and that durability made repurposing easier than replacement. The Fordist history of capacity utilization measurement is therefore not a story of rise and disappearance, but of construction, routinization, recoding, and transition. Brookings made slack visible as a measurable problem. American Keynesianism turned that problem into a functional requirement of stabilization. McGraw--Hill and the Federal Reserve transformed the measure into a durable operational routine. Stagflationary crisis then narrowed its meaning into an inflation sentinel. That sequence matters for what follows. Once the Fordist benchmark had been recoded in this way, the ground was prepared for a new measurement regime in which slack would increasingly be governed not through survey-centered utilization series, but through more portable output-gap and potential-output benchmarks aligned with post-Fordist forms of discipline \citep{MahoneyThelen2009}.


\subsection{The Rise of Econometrics and its Global Sedimentation on Institutions: From Capacity Utilization to Output-Gap Measurement through Post-Fordist Times}
\label{subsec:22_postfordist_capacity_utilization}

% =============================================================================
% §2.2 — The Rise of Econometrics and its Global Sedimentation on Institutions
% R1–R12 Prose Draft (Opus Rewrite)
% Last updated: 2026-03-11
% =============================================================================
% NOTE: Compile with natbib (author-year). All \citep / \citet keys mapped
% against the dissertation master .bib.
%
% MISSING FROM .BIB (must be added before compilation):
%   Foster (2022) — used in R9. Suggested entry:
%     @article{Foster2022,
	%       author  = {Foster, Chase},
	%       title   = {Constructing Consensus: Technocratic Governance and
		%                  the Politics of Monetary Authority after 2008},
	%       journal = {<CONFIRM JOURNAL>},
	%       year    = {2022}
	%     }
%
% DUPLICATE ENTRIES IN .BIB (will cause warnings — deduplicate):
%   - Jessop2002 appears twice
%   - Best2025 appears twice
%   - Sklair2000 appears twice
%   - LofaroMatamorosRochon2025 vs LofaroEtAl2025 (same work, different
%     author spellings: Alessio/Antonio, Gabriel/Gonzalo — verify and keep one)
%
% YEAR DISCREPANCIES (prose year vs .bib year — natbib resolves from .bib):
%   -  & Voigts: prose says 2024, .bib says 2025 → using .bib key
%   - Kavřík: prose says 2026, .bib says 2025 → using .bib key
%   - Dipres: prose says 2024, .bib says 2025 → using .bib key
% =============================================================================


% ---- R1 — Opening Hinge ----

Post-Fordism, in Jessop's regulationist formulation, names not a periodization
label but a reorganization of governance space under conditions of spatiotemporal
compression---the acceleration of capital circulation, the rescaling of state
authority toward supranational venues, and the displacement of nationally bounded
Keynesian routines by portable benchmark architectures legible to transnational
expert fractions and technopol networks
\citep{Jessop2002,Jessop2004,Sklair2000}. The governance demand this compression
creates is historically specific: inherited Fordist routines---national surveys,
dashboard indicators, discretionary fiscal instruments---cannot travel across the
jurisdictions and institutional scales that globalized accumulation now requires.
What emerges in response is a first tendency within post-Fordist macro governance:
the withdrawal of the state from real-time discretionary management of demand and
its replacement by rule-based, pre-announced policy frameworks whose authority
rests not on correct tracking of the cycle but on the institutional consensus that
depoliticized benchmarks can substitute for judgment
\citep{Jessop2002,SaadFilho2018}. Within this first tendency, the output-gap
apparatus operates as a closure device---a portable, institutionally reproducible
governance variable that converts the unruly complexity of slack into a benchmark
legible to surveillance bodies, budget offices, and central banks simultaneously.
It does not track the cycle so much as close the governance problem that
spatiotemporal compression has opened. Yet closure devices of this kind are
themselves spatio-temporal fixes in Jessop's sense: they defer rather than resolve
the contradictions of governance under globalized accumulation, displacing
stagnation risks and crisis pressures across geographies and temporal
horizons---until the conditions underwriting the fix shift, and the apparatus must
repair itself or stretch beyond its original institutional design
\citep{Jessop2004}.


% ---- R2 — Volcker Shock as Pivot ----

The decisive critical juncture through which this first tendency consolidates is
the Volcker Shock, but it should not be narrated as a simple victory of
anti-inflation orthodoxy over Keynesian profligacy. \citet{Konings2007}
demonstrates that the institutional foundations of U.S.\ structural power in
global finance were already being reworked through the re-emergence of
international capital markets from the 1960s onward, and that the monetarist turn
is better read as a reconfiguration of the relationship between U.S.\ state
power, American finance, and the international monetary order than as mere
submission to external discipline. Krippner's long-run definition of
financialization---accumulation increasingly organized through financial channels
rather than commodity production---gives that pivot its structural backdrop
\citep{Krippner2005}, while \citet{Krippner2007} traces how central bank
transparency itself became a governance technology within the neoliberal order.
\citet{MorganSheehan2015} serve as a historiographic corrective against reducing
turning points to abstract theory disputes: elite narratives of reform
characteristically present the reorganization of institutional power as technical
correction, not as a political settlement with distributional consequences. Read
together, these accounts make the Volcker moment legible not as a monetarist
experiment but as the historical hinge through which a new benchmark
regime---one premised on the credibility of depoliticized rules rather than the
discretion of policy authorities---became institutionally plausible.

The theoretical legitimation for this regime shift had been constructed in the
preceding decade through the Lucas critique, which \citet{Johnson2024}
reconstructs not as a laissez-faire argument but as an epistemological claim about
the conditions of possibility for scientific governance. Lucas's program was
premised precisely on the hope that states and scholars could make the economy
visible and controllable---but only through stable, pre-announced policy rules
that eliminated the Knightian uncertainty discretionary intervention produces.
Discretionary policy, by making future governance inscrutable to agents, destroyed
the rational expectations on which econometric legibility depended; it did not
merely fail to stabilize but actively undermined the epistemic infrastructure that
would have made stabilization scientifically defensible. The Volcker shock
operates within this emerging disciplinary consensus as a discretionary act whose
purpose is to establish the credibility that formal rules would later
institutionalize---not a rule-based intervention in itself, but the violent
founding moment of a regime that would subsequently present itself as
rule-governed. What the 1980s achieved through institutionalized credibility among
monetary policy agents, the 1990s would formalize as portable governance
architecture: the Taylor rule, inflation targeting, and central bank independence
statutes converted discretionary consensus into a framework defensible as
technically neutral \citep{Johnson2024,Taylor1993}.


% ---- R3 — Settlement Package Named ----

What the Volcker turn ultimately settles is a durable monetary-policy package that
\citet{SaadFilho2018} names the New Monetary Policy Consensus. In his
formulation, neoliberal monetary policy does not merely supply a superior
stabilization technique; it constructs a political regime around inflation
targeting, central bank independence, and floating exchange rates---a regime
presented as best practice since the late 1980s, whose distributional consequences
are rendered invisible precisely because the framework claims technical
neutrality. \citet{LofaroMatamorosRochon2025} reinforce this diagnosis from a Post-Keynesian
and Sraffian standpoint, showing that the NMPC's governance of interest rates and
inflation targets carries distributive effects that the consensus framework
systematically externalizes. The package matters for this section because it
supplies the closure language through which slack is recoded: no longer
principally idle capacity to be mobilized through demand management, but an output
gap to be monitored, constrained, and interpreted through inflation expectations
and nominal credibility. In other words, the post-Volcker settlement converts
macroeconomic measurement from a demand-management lever into a rule-consistent
benchmark apparatus---one designed not to track the cycle in real time but to
replace real-time discretion with a rule that could be defended as technically
neutral across institutional venues. That is the T1 formalization: the governance
problem opened by spatiotemporal compression is closed not by better measurement
but by measurement recoded as rule.


% ---- R4 — Institutional Comparative Spine ----

This settlement becomes durable because it is embedded in a comparative
institutional spine rather than resting on one central bank model alone. The
apparatus sediments across four privileged venues: the IMF, where potential output
and the output gap operate as surveillance variables, fiscal-position benchmarks,
and policy-recommendation inputs; the OECD, where they serve as medium-term,
inflation-consistent benchmarks for Economic Outlook work and country reviews; the
CBO, where they are folded into budget baselines under explicit medium-term
closure assumptions; and ECFIN, where they are embedded directly in
structural-balance calculations and rule-based fiscal governance through the
Stability and Growth Pact. This is the moment when the output gap ceases to be
one indicator among many and becomes a family of routinized
devices---differentiated by institutional function but unified by a common closure
grammar that recodes slack as an inflation-risk variable amenable to
rule-consistent governance. The post-Fordist rise of the apparatus is therefore a
story not of theoretical convergence alone but of institutional standardization
across distinct policy venues, each of which gives the same benchmark a different
governance lever while preserving its closure architecture. The two focal cases of
this dissertation---the United States and Chile---represent structurally distinct
articulations of this benchmark family: in the United States, the apparatus is
most tightly embedded in budget-baseline interpretation and policy-rule discourse
under financialized post-Fordism; in Chile, the same benchmark arrives through
open-economy transmission logic and structural-balance governance, shaped by
external constraint and the earliest concentrated experience of neoliberal
restructuring in the periphery. Together, the two cases make legible a controlled
comparison of how the same measurement apparatus produces different governance
effects across core and peripheral institutional settings.


% ---- R5 — IMF + OECD Block ----

Within that spine, the IMF and OECD occupy the supranational benchmark tier. At
the IMF, potential output and the output gap function as surveillance variables,
fiscal-position benchmarks, and policy-recommendation inputs---a triple role that
makes the device load-bearing across Article IV consultations, fiscal
sustainability assessments, and program conditionality. The post-GFC
methodological lane has added a layer of self-awareness: revision audits,
uncertainty quantification, country-specific customization, hysteresis
recognition, symmetry-bias diagnostics, and warnings about real-time
misidentification now accompany what was previously an unqualified point estimate
\citep{BarkemaGudmundssonMrkaic2020,CerraSaxena2017,CerraFatasSaxena2023,%
	AiyarVoigts2024}. At the OECD, the device is framed differently---as a
medium-term inflation-consistent benchmark embedded in Economic Outlook
projections, country surveys, structural-balance calculations, and long-run
scenarios---with post-crisis work emphasizing revision control and
anti-cyclicality rather than abandonment of the estimator
\citep{TorresMartin1990}. The result is revealing: post-GFC reflexivity does not
displace the output gap from the apparatus but repackages it as a more cautious
and better-calibrated benchmark, one whose survival depends on absorbing critique
rather than resolving it. That pattern---reflexive adaptation without paradigm
exit---is one of the clearest signatures of how the apparatus reproduces itself
under stress.


% ---- R6 — CBO + ECFIN Block ----

The CBO and ECFIN reveal a different and more coercive embedding of the same
benchmark family. In the U.S.\ case, potential GDP and the output gap are tied
directly to the budget baseline and to standardized interpretations of the budget
balance under explicit medium-term closure assumptions---an architecture that
makes the CBO's output-gap estimate consequential not because it is epistemically
superior but because it is administratively binding
\citep{CBO2004PotentialGDP}. In the EU case, potential growth and output gaps are
folded into rule-based governance through the Stability and Growth Pact, the
European Semester, and structural-balance calculations, where they adjudicate what
counts as fiscal space, discipline, or excess \citep{HavikEtAl2014}. Here the
output gap is no longer only a surveillance device; it becomes part of the
machinery through which fiscal sovereignty is constrained by a measurement routine
that presents itself as technically neutral. The same benchmark family that in one
institutional venue appears as assessment appears in another as
enforcement---and it is precisely this functional versatility that makes the
apparatus durable across governance settings that differ in everything except
their reliance on the output gap as a closure variable.


% ---- R7 — U.S. Focal Case ----

The U.S.\ focal case condenses this post-Fordist settlement with particular
clarity. The technical map identifies the U.S.\ lane as combining CBO closure
logic, IMF methodology repair, policy-rule embedding through the Taylor
tradition, contemporary filter sensitivity, and a retained Blanchard--Quah
structural-decomposition branch in the background
\citep{CBO2004PotentialGDP,AdamsCoe1990,Taylor1993,Alichi2015}. Its defining
asymmetry is that the output gap is especially tightly linked to budget-baseline
interpretation and policy-rule discourse---a linkage that makes the benchmark
consequential not only for monetary but for fiscal governance under financialized
conditions. That makes the U.S.\ case useful not because it is universal but
because it reveals the strongest articulation between financialized macro
governance, benchmark authority, and fiscal interpretation. \citet{Kavrik2026}
sharpens this point by demonstrating that filter-based output-gap estimates in the
U.S.\ policy-rule context exhibit substantial real-time instability---revisions
large enough to alter the sign of the gap and, with it, the policy signal the
benchmark is supposed to deliver. In this configuration, the output gap does not
merely describe slack; it organizes how slack becomes administratively intelligible
for macroeconomic management, even as the real-time reliability of that
organization remains structurally precarious.


% ---- R8 — Chile Focal Case ----

Chile, by contrast, shows how the same apparatus family is articulated differently
under open-small-economy conditions shaped by external constraint, commodity
dependence, and the earliest concentrated experience of neoliberal restructuring
in the periphery. The focal map identifies the Chile lane as combining the IMF
apparatus core, BCCh apparatus core, BCCh real-time and revision governance,
institutional doctrine from BCCh, Hacienda, and Dipres, and an IMF
open-small-economy transmission companion
\citep{MagudMedina2011,FuentesGredigLarrain2008,BCCh2017CrecimientoTendencial,%
	Hacienda2023PIBNoMinero,Dipres2025BCA2024}. Its defining asymmetry---stronger
structural-balance embedding and open-economy transmission logic---gives the
benchmark a different governance lever than in the U.S.\ case: here the output
gap mediates not primarily budget-baseline interpretation but the relationship
between commodity-cycle exposure, inflation pass-through, and fiscal-rule
compliance. \citet{FigueroaForneroGarcia2019} make visible a further layer of
this embedding by documenting how output-gap revisions in the Chilean context
propagate through inflation forecasts and fiscal-balance calculations---a revision
politics in which the benchmark's instability is not merely epistemic but has
direct fiscal consequences. Chile thus operates as the clearest comparative
reminder that portability does not imply homogeneity: the benchmark survives
across institutional settings, but its governance leverage depends on the
macro-structural and institutional terrain in which it is embedded.


% ---- R9 — Opening the Post-GFC Break ----

The crisis of 2008 does not dislodge this settlement; it exposes the conditions on
which the first tendency's viability has always depended. \citet{Best2025}
reconstructs the political economy of monetary depoliticization in the UK to show
that rule-based governance frameworks hold only while the narrow technical devices
on which they rest retain credibility---and that credibility requires a
historically specific set of macroeconomic conditions: moderate shocks, anchored
expectations, and institutional environments where pre-announced targets can
absorb short-run deviations without triggering the political backlash
depoliticization was designed to foreclose. Post-2008---and more sharply after the
COVID disruption---these conditions erode, and a second tendency becomes visible:
the statistical and information infrastructure assembled during three decades of
benchmark governance now makes real-time intervention technically feasible in ways
that were unavailable to Fordist or early post-Fordist states
\citep{Foster2022}. Central banks do not merely repair the old rule-based
architecture; they begin to reoccupy the discretionary governance space that the
first tendency had explicitly vacated---expanding mandates, deploying
unconventional instruments, and interpreting their targets with a flexibility that
the original depoliticized framework was built to prohibit. The filtering lane
reaches a parallel conclusion from the methodological side: output-gap estimation
rests on competing devices with different closure assumptions, revision
properties, and crisis sensitivities, and there is no single technical
winner---only a portfolio of tools whose coexistence is itself evidence that the
benchmark can no longer be defended as a settled, rule-consistent variable
\citep{Hamilton2018,QuastWolters2022,DurandFornero2024}. The post-GFC period
should therefore be read not as collapse but as the moment when T1 institutional
architecture begins to bear T2 operational demands---a mismatch whose sharpest
expression is not methodological plurality but the legitimacy crisis it produces.


% ---- R10 — Legitimacy Crisis of Depoliticization ----

The ECB provides the sharpest meso-institutional illustration of this T1/T2
mismatch. \citet{ArgyroulisVagdoutis2025} show that the ECB's early strategy
statements were built for a T1 world: tightly defined performance criteria, a
narrow price-stability mandate, and an institutional self-understanding premised
on the claim that technocratic governance required neither democratic
authorization nor discretionary judgment beyond the application of pre-announced
rules. The 2021 strategy revision installs a different operational reality. It
broadens the framework, expands discretionary space, incorporates climate and
distributional considerations, and reinterprets the mandate in ways that
Argyroulis and Vagdoutis characterize as rendering it ``broad and vague rather
than clear and narrow'' (p.~14)---a formulation that captures precisely the
institutional diagnosis at stake. What they name is not a failure of
implementation but a structural contradiction: the apparatus must stretch its
discretion to absorb post-crisis demands---multiple objectives, flexible
medium-term orientation, unconventional instruments---yet in doing so it
dismantles the very claim to neutral, rule-bound expertise on which its
depoliticized authority was constructed. The result is not the end of technocratic
governance but a structurally weakened version of it---one that must now justify
its discretion rather than simply assert its technical neutrality, and that can no
longer present the output gap as a self-evidently rule-consistent variable. This
is not a contingent difficulty at the edge of the system; it is the concentrated
institutional expression of a regime whose first-tendency architecture was not
designed to carry the governance demands its own crisis management now requires.


% ---- R11 — Filtering and Methodological Plurality as Political Symptom ----

Once seen from this angle, methodological plurality is not a technical side note
but a political symptom of an apparatus repairing itself for second-tendency
demands using first-tendency tools. The filtering branch---Beveridge--Nelson
decomposition, Hodrick--Prescott, band-pass variants, Hamilton-type structural
critiques, and the later repair literature---does not converge on one universally
superior output-gap estimator; these devices coexist as a portfolio whose
usefulness depends on data-generating structure, intended policy use,
institutional benchmark requirements, and tolerance for revisions
\citep{BeveridgeNelson1981,HodrickPrescott1997,BaxterKing1999,%
	ChristianoFitzgerald2003,Hamilton2018,FrankeKukackaSacht2024,Siemers2025}. No
single technique wins because the governance problem the apparatus must solve has
shifted: what was once a question of selecting the correct filter for a
rule-consistent benchmark is now a question of which combination of devices can
sustain benchmark authority under conditions where the rules themselves are being
renegotiated. The technical branch is therefore saturated for the purposes of
\S2.2---not because one method has been vindicated, but because the apparatus
survives by combining and repurposing coexisting devices while preserving the
output gap as the dominant governance variable. Plurality is contained inside the
benchmark regime rather than opening it to paradigmatic challenge.


% ---- R12 — Exit Ramp to §2.3 and Shaikh ----

The payoff for the rest of the chapter is now legible. Post-Fordist macro
governance globalized the measurement of slack by embedding the output gap in
surveillance, budgeting, structural-balance rules, and central-bank strategy
statements, while financialization and technocratic statecraft made this benchmark
increasingly authoritative---portable enough to travel from Washington to
Santiago, from Article IV consultations to structural-balance calculations, from
Taylor-rule calibrations to ECB mandate revisions. But the very durability of the
apparatus exposed its limits: it is revision-prone, closure-dependent, and
institutionally selective, and its post-crisis survival increasingly depends on
the discretionary repair and broadened mandates that its original depoliticized
architecture was designed to render unnecessary. That is why the heterodox
reopening of capacity utilization in the next section becomes historically
intelligible. It does not arise in an empty theoretical field; it emerges against
the dominance of an output-gap apparatus that has become globally portable
precisely by turning one recoded form of slack into a governing norm---and against
the growing visibility of the contradictions that portability produces when the
governance demands outrun the institutional design.





% ============================================================
\subsection{The Capacity Utilization Controversies in Political Economy Debates}
\label{sec:2.3}
% ============================================================


The debate over the measurement of productive capacity and potential output is not confined to technical economics. When estimates of the output gap serve as the operative benchmark for structural-balance fiscal rules (as they did across the Eurozone in the aftermath of the 2008 crisis), their systematic bias carries observable distributional consequences. \citet{Klar2013} documents how unreliable potential-output and NAIRU estimates constrained fiscal policy in Europe, producing contractionary impulses in economies operating well below productive capacity. \citet{Heimberger2017} show how the European Commission’s model for estimating ‘potential output’ \citep{HavikEtAl2014}, are used to calculate ‘structural budget balances’, which in turn translate into country-specific  structural-balance governance reinforcing pro-cyclical austerity logic. \citet{Heimberger2020} tracks how European Commission
potential-output revisions reshaped fiscal space under crisis conditions, compressing public investment precisely when demand shortfalls were largest. From a classical-keynesian perspective, \citet{Palumbo2015} provides a theoretical restatement of Okun's Law explaining what these findings imply: potential output is endogenous and path-dependent rather than an exogenous supply ceiling. Hence, the measurement apparatus was not merely imprecise but
conceptually misconceived as a governance instrument. Mismeasurement, in this register, is not a technical nuisance---it allocates the costs of under-utilization to workers and the public sectors while insulating financial creditors from adjustment. That the dominant measurement instrument produces such consequences with such regularity raises an immediate question for political economy as a research program: if the dominant measurement instrument is consequentially wrong, what alternative does the tradition offer?

The answer reveals a problem unresolved within Political Economy debates, which expresses itself with the standard practice of research to engage capacity utilization taking the Federal Reserve Board series as an empirical
object taken at face value without further scrutiny on its methodological construction. As showcased above, the FRB capacity utilization series was not produced in a vacuum, but produced as a measurement of governance under a historically contingent institutionalization where: monitoring manufacturing slack for monetary policy purposes within the institutional landscape described in \ref{subsec:21_fordist_capacity_utilization}. Its migration into political economy research happened without systematic evaluation of its adequacy for subjecting theoretical claims to empirical scrutiny. The series existed, covered the relevant variables, carried official
authority, and offered long historical coverage. These are properties that make an instrument useful for governance dashboards. They are not properties that make it adequate for adjudicating theoretical claims about the long-run dynamics of capitalist accumulation; no such distinction was drawn at the point of adoption. The instrument was
recruited from a governance context into a scientific one without the transit being examined.

The historical timing of this migration compounds the problem. The political economy debate over the normal rate of capacity utilization gains its contemporary momentum in the aftermath of the Global Financial Recession in which Political Economy approaches to macroeconomics have gained traction amidst the crisis of the Post-Fordist regime. As portrayed in \ref{subsec:22_postfordist_capacity_utilization}, the output-gap measurements constituted themselves as an statistical apparatus embedded across supranational institutions building the operational infrastructures of fiscal surveillance.  The same class of instrument is simultaneously institutionalized as governance infrastructure by mainstream bodies and recruited as scientific evidence by political economy researchers pursuing opposed theoretical agendas. A further scope constraint operates silently throughout: FRB is an industrial manufacturing survey. In the Fordist era this covered a dominant sector of accumulation. In the Post-Fordist economy it traces a contracting share of output, employment, and investment while being used to characterize economy-wide capacity dynamics. The representational problem deepens across the period of the instrument's heaviest use in political economy research.

The controversy over capacity utilization's long-run behavior carries substantive theoretical stakes that deserve clear statement before their empirical impasse is diagnosed. The Neo-Kaleckian tradition argues that
the normal rate of capacity utilization is endogenous to demand. In this view, firms revise their desired utilization in response to persistent demand conditions: through gradual adjustment of expectations \citep{Lavoie1995}, through micro-founded shift-choice mechanisms in which indivisibility makes utilization respond when capital cannot be smoothly resized \citep{Nikiforos2013}, or through hysteresis in which the history of demand shocks durably reshapes the utilization target itself \citep{Setterfield2020}. The implication is that actual utilization need not converge to any structurally given benchmark: demand conditions shape long-run accumulation trajectories and distributive outcomes are path-dependent. The Sraffian-Keynesian tradition advances the opposed claim. In the classical formulation, the
normal rate is determined by cost structures and competitive conditions, and actual utilization gravitates toward it over the long run, while the speed of convergence might be mediated due sluggish investment and a desired rate with slack to uncertainity of demand surges \citep{Ciccone1986}. 

The contemporary version is more precise: within the Sraffian supermultiplier framework, autonomous demand drives the growth trend of output while investment adjusts to bring capacity in line with demand, so that utilization converges toward a normal rate determined by structural conditions rather than by demand history \citep{Girardi2016}. The convergence claim is buttressed by cross-country evidence: \citet{GahnGonzalez2022} document stationarity of utilization across twenty-one countries and find positive short-run but negligible long-run correlation between utilization and growth---a pattern they read as more compatible with convergence than with permanent Neo-Kaleckian drift.

If utilization gravitates to a structurally given normal rate, then demand-side interventions have only transitory effects on accumulation---the system self-corrects through investment adjustment. If utilization is path-dependent, then persistent demand shortfalls generate durable consequences for both investment and employment, which means macroeconomic governance can reshape accumulation trajectories rather than merely smoothing cycles. Resolving the controversy empirically requires a measurement instrument whose properties carry genuine information about the long-run relationship between actual and normal utilization. 

Nevertheless, that requirement has not been met. Both sides of the controversy require a measure whose stochastic properties (stationary or trending, mean-reverting or drifting) carry information about whether actual utilization converges to a normal rate over the long run. The Federal Reserve Board series cannot provide this, because its statistical behavior is not a property of the underlying economic process but an artifact of construction. Capacity Utilization series built by the FRB combines survey responses on current operating
rates with regression-based capacity estimates that are smoothed against capital stock and time trends, explicitly adjusted to move the resulting series away from an engineering concept of maximum output toward an economic concept of sustainable output under normal operating conditions \citep{Shapiro1989}. The stated design objective is that
the resulting measure has as a byproduct an engineering mean-reversion behavior by construction, and is impeded of registring evidence of chronic departure from normal capacity utilization: 

Testing whether the FRB series is stationary is therefore not a test of whether the actual utilization rate converges to a normal rate fixed by cost conditions; it is a test of whether the Federal Reserve's construction procedure successfully implemented its own design objective. \citet{Nikiforos2016,Nikiforos2021} provide the most
systematic articulation of this problem, arguing that FRB is closer to a cyclical slack indicator than to a measure of long-run movements in normal utilization. The relevant time horizon for the controversy---on the order of the useful life of capital, roughly thirty years---far exceeds the frequency at which the series is informative.

The construction opacity compounds the difficulty. The raw micro-data underlying the series is not publicly available, the weighting procedures combining different survey sources are not fully disclosed,
and the splicing across institutional transitions lacks transparent documentation \citep{Board2017}. The instrument cannot be subjected to independent methodological scrutiny by the researchers who rely on it. The category error is therefore precise: FRB was built to produce a stable, institutionally credible signal for monetary policy governance---stability is a virtue for that purpose and a disqualifying defect for adjudicating a controversy that turns on whether utilization is stable or trending. But even if FRB were replaced by a transparent and well-constructed survey, a deeper problem would remain, one rooted not in the instrument's opacity but in the micro-level premise that all survey-based utilization measures share.

Even if a transparent and well-constructed survey replaced FRB, the deeper problem is that both the Neo-Kaleckian desired rate and the Sraffian normal rate presuppose that firms have a coherent utilization
target, whether endogenously revised or exogenously anchored, that can in principle be elicited through survey instruments. This shared micro-level premise is empirically unsupported. The difficulty begins with scope: FRB covers manufacturing, which in the Fordist era traced a dominant share of accumulation but in the Post-Fordist economy
represents a contracting fraction of output, employment, and investment  \citep{Haluska2023}. The measurement object thus narrows as the theoretical ambition widens. But the problem runs deeper than sectoral coverage. Management and operations research literature converges on a portrait of capacity utilization as a neglected variable at firm level,
managed tactically for short-run performance rather than optimized toward a long-run target rate \citep{Anderson2001,Singh2022}. Planned utilization levels are simultaneously conditioned by competitive
pressures, demand expectations, and cost structures---outcomes of contradictory organizational logics rather than expressions of a unified firm objective \citep{Lieberman1989}.

The structural source of this incoherence is the transformation of corporate governance itself. \citet{vidal2022} provides the most pointed formulation: the Post-Fordist firm is not a unified decision-making agent. Plant-level managers maximize throughput against cost-cutting constraints. Corporate executives manage earnings per
share, share buybacks, and stock price performance under the shareholder-value regime that restructured Anglo-American corporate governance from the 1980s onward \citep{Lazonick2000}. Directors mediate agency conflicts between management and ownership whose objectives are misaligned by design \citep{Fligstein2015}. Even if
every tier responds honestly to a utilization survey, the aggregated response does not represent a coherent firm-level desired or normal rate---it represents the resultant of an internal contest whose outcome
has no clean theoretical interpretation on either side of the controversy. The irony is temporal: shareholder-value capitalism decouples plant-level utilization decisions from corporate accumulation strategy during the period when the political economy measurement debate intensifies. The measurement object becomes less theoretically interpretable exactly when it is most heavily used. The works by Nikiforos recognizes this cluster of problems more clearly than any other participant in the debate, and his attempt to address it through measurement rescue is the most self-reflexive move within the tradition.

\citet{Nikiforos2016} engages the measurement debate by subjecting the FRB series to systematic scrutiny
on its own terms. His critique of construction opacity, respondent ambiguity, and engineered mean reversion is well-founded and independently motivated: the problems identified in the preceding paragraphs are substantially his diagnosis. His proposed rescue is correspondingly analogous: the average workweek of capital, developed by
\citet{Foss1981,Foss1985} among others, is a direct physical measure of machinery use, sidesteps both the construction artifacts of FRB and the micro-level premise problem, since it measures capital intensity rather than eliciting managerial assessments of a desired operating rate. \citet{Nikiforos2021} deepens the case by arguing that AWW and related physical proxies are conceptually better matched to the long-run horizon relevant to the controversy. This is a genuine methodological advance. The question is whether it can settle the dispute. AWW trends where FRB does not, and the cointegration of FRB on AWW is read as evidence that FRB's stationarity reflects measurement error rather than
true convergence. But this leaves open whether the proxy's non-stationarity is evidence for endogenous normal utilization or simply a property of a different measurement object---and the single-equation ARDL framework within which the evidence is produced cannot adjudicate that question from within.

\citet{GahnGonzalez2020} press the empirical side of this concern, arguing that the ARDL evidence does not robustly establish that AWW delivers the stochastic-trend story the Neo-Kaleckian interpretation requires.
Their critique is effective as a rebuttal of Nikiforos's specific econometric exercise. But it does not itself escape the same methodological terrain. The Sraffian-convergence response to the controversy relies on a set of stationarity tests, proxy admissibility arguments, and transitory-shock evidence \citep{GahnGonzalez2022,Deleidi2024}; yet demonstrating that the Neo-Kaleckian case is empirically insufficient is not the same as positively identifying the long-run adjustment mechanism through which utilization is supposed to converge. What the debate has established, then, is not which tradition is right but that the instruments available to either side cannot settle the
question. Both sides of the controversy appear to remain confined to single-series proxy disputes and single-equation testing strategies that may not be well suited to recover the structural object at stake. If the long-run relationship between actual utilization and the capacity benchmark is shaped by the joint dynamics of accumulation, distribution, and technical change, it may be better understood as a system property than as a property of any single observable series. Unit root tests, proxy comparisons, and conditional error-correction models applied to individual series are not well suited to identify it. Whether the controversy is blocked by insufficient data or by the limits of the identification strategies available to it remains an open question---but one that the measurement debate in its current form does not appear equipped to resolve.


% ============================================================
% §2.4 Shaikh's Critique and the Requirement for Structural Identification
% Chapter 2 — Historical Trace of Capacity Utilization Measurements
% Redraft: three-paragraph version | Session 5 | 2026-03-10
% Per audit: EMU table → Appendix~\ref{app:emu_lineage}
%            §3 owns ARDL/ECM exposition — one sentence here only
%            Dadda2003 dropped (unverified bib key)
% ============================================================

\subsection{Shaikh's Critique and the Requirement for Structural Identification}
\label{sec:24_shaikh_structuralid}

% ---------------------------------------------------------------
% P1 — Shaikh as device-interrogator + physical measures lineage
%       + two-level FRB critique + Nikiforos formalization
%       + cross-reference to EMU appendix table
% ---------------------------------------------------------------

Shaikh's interrogation of the measurement apparatus is the only intervention
in the political economy literature that proceeds independently of the
utilization controversy itself --- not as a contribution to the
Kaleckian--Sraffian debate but as a critical assessment of the data
infrastructure underlying the rate of profit. The tradition from which it
emerges begins with \citet{Foss1963}, who constructed the Electric Motor
Utilization index by comparing actual electricity consumption in
manufacturing with the theoretical maximum implied by installed motor
horsepower, producing a physically grounded intensity measure that registered
higher utilization rates and greater long-run variation than the Federal
Reserve series. \citet{Shaikh1992} extended the method, building a
continuous EMU series from 1899 to 1963 using Census data and adapting
\citeauthor{Christensen1969}'s (\citeyear{Christensen1969}) estimates, then
splicing it with McGraw-Hill survey data to construct a utilization measure
covering 1947 to 1985; the methodology underlying this lineage is
documented in Appendix~\ref{app:emu_lineage}. Shaikh's critique of the FRB
series operates on two levels. 

At the institutional one, the Federal
Reserve combines survey data through undisclosed procedures and adjusts the
resulting estimates with the stated objective that the measure should not be
chronically below normal, making stationarity a design feature rather than a
discovered property of the data \citep[pp.~187--188]{Shapiro1989}. At the
conceptual level, physical measures reveal non-stationary, asymmetric
fluctuations with greater amplitude than FRB, suggesting that the
survey-based series suppresses precisely the long-run patterns that matter
for the theoretical question. \citet{Nikiforos2016} later formalized this
critique through the plant-manager thought experiment: because adding or
dropping shifts changes the definition of normal operating conditions, the
survey denominator moves with the numerator, making the instrument
self-referential with respect to the object it purports to measure.

% ---------------------------------------------------------------
% P2 — Shaikh's advance (unbalanced growth) + architectural limit
%       (single-equation assumes away distributive channel)
%       §3 owns the ARDL/ECM exposition — one sentence only here
% ---------------------------------------------------------------

Shaikh's proposed solution --- recovering productive capacity as the
long-run component of the output--capital cointegrating relation under a
maintained normal-utilization closure --- represents an advance over both
survey-based measures and proxy comparisons; the method and its econometric
implementation are examined in \S\ref{sec:ch3}. The substantive empirical
finding that advance produces is unbalanced growth: the growth rate of
productive capacity persistently diverges from the growth rate of capital
accumulation, which breaks the balanced-growth assumption shared by both the
Kaleckian and Sraffian positions, where the transformation elasticity
$\theta$ linking capacity expansion to capital accumulation is implicitly
treated as unity. Shaikh showed it is not. The limit of the approach,
however, is architectural. The single-equation specification captures
unbalanced growth through a fixed transformation elasticity, when recent
theoretical work re-interpreting \citet{Kaldor1957} suggests that the
elasticity is distributionally conditioned: following \citet{Cajas2024},
distribution drives mechanization intensity, which drives the rate at which
capital accumulation translates into capacity expansion, making the
transformation elasticity $\theta(e;\,\bar{\Lambda})$ regime-specific rather
than structurally constant. A bivariate specification cannot recover the
transformation elasticity as a system property because it assumes away the
distributive channel through which mechanization operates --- and the
discretionary dummy structures it requires to sustain any long-run
relationship signal precisely that distribution matters, without providing a
variable through which it can act.

% ---------------------------------------------------------------
% P3 — Forward link to critical replication
%       Trimmed: dropped "What does adequate identification require?"
%       rhetorical opener; kept three-stage architecture statement
% ---------------------------------------------------------------

This reframing defines the methodological requirement: a cointegration
framework that endogenizes the transformation elasticity, captures the joint
dynamics of accumulation, distribution, and capacity transformation, and
recovers utilization as bounded movement within an institutionally determined
space rather than as gravitation toward a fixed benchmark. The empirical
question is whether Shaikh's measurement pipeline, when subjected to critical
replication and extended to a system framework, bears out this confined
structure --- and where its identification breaks.






\section{Shaikh's Framework for Measuring Capacity Utilization}
\label{sec:shaikh_framework}

Shaikh's utilization measure is organized around a constructed benchmark
level of output, $Y_t^p$, interpreted as productive capacity:
%
\begin{equation}
	u_t \equiv \frac{Y_t}{Y_t^p},
	\qquad
	\ln u_t = \ln Y_t - \ln Y_t^p.
\end{equation}
%
The identifying convention is a maintained long-run closure in which
utilization converges to its normal level ($u = 1$). Under this closure,
productive capacity is not separately observed; it is inferred as the
long-run component of the output--capital relation, and utilization is
recovered residually as the deviation from that component. Shaikh
operationalizes this inference using an ARDL model in error-correction form
within the \citet{Pesaran2001} bounds-testing framework. The replication
that follows saturates this identification device rather than replacing it.


% ============================================================
% §3.1 THE MEASUREMENT APPARATUS: GPIM
% ============================================================

\subsection{The Measurement Apparatus: Shaikh's Generalized Perpetual
Inventory Method}
\label{subsec:gpim}

The capital stock series entering the identification are not raw BEA
aggregates. They are constructed through Shaikh's Generalized Perpetual
Inventory Method (GPIM), which closes a gap in the national accounts: the
standard micro-level perpetual inventory identity does not survive
aggregation under the Fisher-ideal chain-weighting that the BEA uses to
construct aggregate capital stocks.\footnote{The data objects are drawn from
Appendix~6.8 of \citet{Shaikh2016}, accessible at the companion website
\texttt{realeconomicanalysis.com}. The source spreadsheet is
\texttt{\_Appendix6.8DataTablesCorrected.xlsx}. A data audit verifying
series alignment is documented in script \texttt{90\_data\_audit.R}
(output in \texttt{output/CriticalReplication/data\_audit/}).}

The GPIM accumulation rule for the aggregate current-cost capital stock is:
%
\begin{equation}
	\label{eq:gpim_accumulation}
	K_t \;=\; IG_t \;+\; z_t^{\ast} \cdot K_{t-1},
	\qquad
	z_t^{\ast} \;\equiv\; (1 - z_t)\,\frac{p_t^K}{p_{t-1}^K},
\end{equation}
%
where $IG_t$ is aggregate gross investment in current prices, $z_t$ is the
aggregate depletion rate, and $z_t^{\ast}$ compresses two conceptually
distinct operations into a single coefficient: the fraction $(1 - z_t)$ of
the stock that physically survives the period, and the revaluation of
that surviving stock from period-$(t-1)$ to period-$t$ prices. Survival is
physical and comes first; revaluation is the accounting price adjustment
applied to whatever survived. The two operations commute algebraically, but
conceptually the distinction matters: the gross stock counts all surviving
machines equally (depletion governs exit), while the net stock discounts
them by accumulated value attrition (depreciation governs book value).
Productive capacity depends on the physical quantity of installed capital,
not on its accounting valuation --- a machine 80\% depreciated still
produces at full service capacity until it retires. The gross stock is
therefore the theoretically appropriate capacity measure; the net stock
confounds a price concept with a physical concept.

For the US corporate sector, Shaikh applies three adjustments to the BEA
official series: substitution of 1993-vintage BEA depletion rates for the
post-revision geometric rates, pre-war initial values from IRS balance-sheet
data rather than BEA backward extrapolation, and a Great Depression
write-down correction that reduces the 1947 corporate net stock by
approximately 28\% relative to the BEA 2011 official figure. The resulting
series --- $\text{KGCcorp}_t$ (gross, GPIM-adjusted) and
$\text{KNCcorp}_t$ (net, GPIM-adjusted) --- together with the BEA benchmark
$\text{KNCcorpbea}_t$, enter the replication as the capital-stock objects.
Full construction formulas and the GPIM derivation from the micro-level
identity through the aggregation problem to the corporate adjustments are
provided in Appendix~\ref{app:gpim}.


% ============================================================
% §3.2 DATA CONSTRUCTION: COMMON DEFLATOR + NET-GROSS PAIRING
% ============================================================

\subsection{Data Construction: the Common Deflator and the Net--Gross
Pairing}
\label{subsec:data_pairing}

% --- Common deflator ---

Both output and capital are deflated by a single price index --- the GDP
implicit price deflator $p_t$ (BEA NIPA Table~1.1.9, line~1) --- yielding
$y_t = \ln(\text{VAcorp}_t / p_t)$ and
$k_t = \ln(\text{KGCcorp}_t / p_t)$. Shaikh's identification requires a
common numeraire: separate deflators would introduce a spurious
relative-price wedge into the output--capital ratio, violating the
profit-rate consistency constraint derived in \citet[Appendix~6.6]{Shaikh2016}.
One tension internal to this specification must be noted. The
common-deflator real stock $\tilde{K}_t = K_t / p_t$ does not obey the
GPIM real accumulation rule, which governs $K_t^R = K_t / p_t^K$. The two
objects differ by the relative price of capital goods:
$\tilde{K}_t = (p_t^K / p_t) \cdot K_t^R$. The common deflator preserves
ratio consistency at the cost of breaking the stock's own law of motion ---
a confinement limit of Shaikh's data environment, not an invalidation of
the identification.

% --- Net-gross pairing: error-term derivation ---

The capacity closure defines productive capacity with respect to the gross
stock:
%
\begin{equation}
	\label{eq:capacity_closure}
	Y_t^p \;=\; A_0\, e^{bt}\,
	\bigl(\tilde{K}_t^{G}\bigr)^{\theta(\Lambda)},
	\qquad
	\ln u_t \sim I(0).
\end{equation}
%
All four logically possible output--capital pairings recover the same
elasticity $\theta$ under this closure, but they differ in residual
composition. The output identity $\text{NVA}_t = \text{GVA}_t(1 - s_t^D)$
and the stock identity
$\ln \tilde{K}_t^N = \ln \tilde{K}_t^G + \ln \psi_t$ generate each
pairing's error term by substitution:
%
\begin{equation}
	\label{eq:four_pairings_errors}
	\varepsilon_t^{GG} = \ln u_t, \qquad
	\varepsilon_t^{NG} = \ln u_t + \ln(1 - s_t^D), \qquad
	\varepsilon_t^{NN} = \ln u_t + \ln(1 - s_t^D) - \theta\ln\psi_t, \qquad
	\varepsilon_t^{GN} = \ln u_t - \theta\ln\psi_t,
\end{equation}
%
where $s_t^D \equiv D_t / \text{GVA}_t$ is the depreciation share and
$\psi_t \equiv K_t^N / K_t^G$ is the net-to-gross ratio. Cointegration
requires each residual to be $I(0)$. Since $s_t^D$ trends upward (secular
rise in CFC/GVA under post-2013 BEA data releases) and $\psi_t$ trends
downward (shift toward shorter-lived assets), the pairings that load these
trending variables into the residual require progressively stronger
auxiliary assumptions: GG requires none beyond the closure; NG requires
$s_t^D$ trend-stationary; NN requires both $s_t^D$ and $\psi_t$
trend-stationary; GN requires $\psi_t$ trend-stationary. The full
derivation with proofs is in Appendix~\ref{app:four_pairings}.

The theoretically grounded pairing is NG: net value added on gross capital
stock. The capacity closure is defined on the gross stock (physical
survival, not book depreciation). The GPIM stock-flow consistency rules
entail cointegration of NVA --- not GVA --- with the gross stock under the
common deflator. And VAcorp (Shaikh's corrected corporate net value added)
is the series published in Shaikh's online supplement --- the choice is
within Shaikh's data perimeter, not an external substitution. The ratio
diagnostic confirms alignment:
$\text{VAcorp}/\text{KGCcorp} \approx R_{\text{corp}}$, Shaikh's canonical
corporate profit rate.

A fourth variable, the corrected corporate profit share
$\text{Profshcorp}_t$, is constructed in the same data environment as the
empirical proxy for the rate of exploitation $\ln e_t$ entering the $m = 3$
VECM system. It does not enter the single-equation identification.
The series entering all estimation stages are therefore
$y_t$, $k_t$, and (for the trivariate system) $\ln e_t$, with $T = 65$
annual observations over 1947--2011.


% ============================================================
% §3.3 S0: FAITHFUL REPLICATION
% ============================================================

\subsection{Faithful Replication of the Capacity Benchmark}
\label{subsec:S0}

% --- What must be recovered ---

Shaikh's published capacity benchmark is a single point in a large
specification space. Before that point can be interrogated, it must be
recovered. The output--capital relation is estimated in ARDL
error-correction form:
%
\begin{equation}
	\label{eq:S0_ecm}
	\Delta y_t
	=
	\alpha
	+ \sum_{i=1}^{p-1} \psi_i\,\Delta y_{t-i}
	+ \sum_{j=0}^{q-1} \varphi_j\,\Delta k_{t-j}
	+ \phi_1 y_{t-1}
	+ \phi_2 k_{t-1}
	+ \gamma t
	+ \sum_h \delta_h DM_{h,t}
	+ \varepsilon_t,
\end{equation}
%
estimated by OLS via the \texttt{ARDL} package for \textsf{R}
\citep{Natsiopoulos2024}. The benchmark specification is fixed at
$m_0 = (2,\, 4,\, \text{III},\, \{d56, d74, d80\})$: lag orders ARDL(2,4),
deterministic Case~III (unrestricted intercept, no trend in the long-run
relation), and Shaikh's reported dummy structure. Admissibility is assessed
at three levels: the bounds $F$-test on the joint null
%
\begin{equation}
	\label{eq:S0_admissibility}
	H_0:\; \phi_1 = \phi_2 = 0
\end{equation}
%
is the primary gate; the bounds $t$-test on the error-correction term
provides a robustness diagnostic where case-specific critical values are
defined (Cases~III and V); and RMSE between $\widehat{u}_t$ and Shaikh's
published series assesses replication fidelity.

The bounds framework distinguishes five deterministic cases according to
whether intercept and trend are excluded, restricted to the long-run
relation, or left unrestricted in the short-run dynamics.
Table~\ref{tab:pss_cases_yp} reports the implied long-run benchmark for
each case. Shaikh's benchmark specification corresponds to
Case~III~($\star$).

\begin{table}[!ht]
	\centering
	\small
	\caption{Deterministic cases in the ARDL bounds framework}
	\label{tab:pss_cases_yp}
	\begin{tabular}{p{0.8cm}p{3.6cm}p{4.2cm}p{2.0cm}p{2.5cm}}
		\hline\hline
		Case & Deterministic treatment
		& Long-run benchmark $\widehat{y_t^p}$
		& Admissibility gate
		& Robustness diagnostic \\
		\hline
		I   & No intercept, no trend
		& $\widehat{\theta}\,k_t$
		& $F$-bounds & --- \\[2pt]
		II  & Restricted intercept, no trend
		& $\widehat{\mu} + \widehat{\theta}\,k_t$
		& $F$-bounds & not defined \\[2pt]
		III$\star$
		& Unrestricted intercept, no trend
		& $\widehat{\mu} + \widehat{\theta}\,k_t$
		& $F$-bounds & $t$-bounds \\[2pt]
		IV  & Unrestricted intercept, restricted trend
		& $\widehat{\mu} + \widehat{\theta}\,k_t + \widehat{\delta}\,t$
		& $F$-bounds & not defined \\[2pt]
		V   & Unrestricted intercept, unrestricted trend
		& $\widehat{\mu} + \widehat{\theta}\,k_t + \widehat{\delta}\,t$
		& $F$-bounds & $t$-bounds \\
		\hline
	\end{tabular}
	\par\vspace{0.3em}
	\footnotesize $\star$ Shaikh's benchmark case: ARDL(2,4), Case~III,
	dummies $\{d56, d74, d80\}$.
\end{table}

Conditional on admissibility, the long-run multipliers are recovered via the
standard ARDL reparameterization:
%
\begin{equation}
	\label{eq:S0_lrmult}
	\hat{a} = \frac{\hat{\alpha}}{1 - \sum_i \hat{\psi}_i},
	\qquad
	\hat{\theta} = \frac{\sum_j \hat{\varphi}_j}{1 - \sum_i \hat{\psi}_i},
	\qquad
	\hat{c}_h = \frac{\hat{\delta}_h}{1 - \sum_i \hat{\psi}_i},
\end{equation}
%
and the capacity benchmark and utilization series are:
%
\begin{equation}
	\label{eq:S0_yp}
	\widehat{y_t^p}
	=
	\hat{a} + \hat{\theta}\,k_t + \sum_h \hat{c}_h\,DM_{h,t},
	\qquad
	\widehat{u}_t
	\;\equiv\;
	\exp\!\left( y_t - \widehat{y_t^p} \right).
\end{equation}
%
The parameter $\hat{\theta}$ is therefore not a direct structural object but
a ratio whose precision depends on the distance of
$1 - \sum\hat{\psi}_i$ from zero. Shaikh's published benchmark values are
$\hat{\theta} = 0.6609$, $\hat{a} = 2.1782$, and long-run dummy multipliers
$\hat{c}_{d74} = -0.8548$, $\hat{c}_{d56} = -0.7428$,
$\hat{c}_{d80} = -0.4780$.


% --- S0 Results ---

\subsubsection{Results}
\label{subsec:S0_results}

Table~\ref{tab:S0_bounds} reports the bounds test under the five PSS cases.
Only Case~II passes the F-test at $\alpha = 0.05$ ($F = 4.16$,
$p = 0.049$); the remaining four cases fail to reject the null of no
cointegration. Case~II does not admit the t-test (restricted intercept
entails a circular null), so the admissibility verdict rests on the F-test
alone. Under Cases~II--III (no trend), $\hat{\theta}_0 = 0.745$
($\text{SE} = 0.147$); under Cases~IV--V (with trend),
$\hat{\theta}_0 = 0.541$ ($\text{SE} = 0.682$) while the trend coefficient
is insignificant ($p = 0.756$). All three step dummies are individually
insignificant across every case ($p > 0.49$).

\begin{table}[!ht]
	\centering
	\small
	\caption{S0 bounds test results: ARDL(2,4) under five PSS cases}
	\label{tab:S0_bounds}
	\begin{tabular}{cccccccc}
		\hline\hline
		Case & $F$-stat & $p_F$ & $t$-stat & $p_t$
		     & $\hat{\theta}$ & $\hat{\alpha}$ & Pass? \\
		\hline
		I    & ---   & ---   & ---    & ---   & ---   & ---     & No  \\
		II   & 4.16  & 0.049 & ---    & ---   & 0.745 & $-$0.12 & Yes \\
		III  & 1.24  & 0.758 & $-$1.33 & 0.714 & 0.745 & $-$0.12 & No  \\
		IV   & 0.84  & 0.982 & ---    & ---   & 0.541 & $-$0.13 & No  \\
		V    & 0.91  & 0.953 & $-$1.35 & 0.894 & 0.541 & $-$0.13 & No  \\
		\hline
	\end{tabular}
	\par\vspace{0.3em}
	\footnotesize
	Bounds tests at $\alpha = 0.05$, asymptotic critical values ($T = 1000$).
	$\hat{\theta}$: long-run output--capital multiplier.
	$\hat{\alpha}$: error-correction speed.
	The t-test is structurally undefined for Cases~II and IV
	(restricted deterministic; see text).
	Shaikh's published $d = 0.661$ (pre-2013 BEA data releases).
\end{table}

The 27\% shift in $\hat{\theta}$ upon trend inclusion could reflect the
particular output--capital pairing rather than a property of the method.
Table~\ref{tab:S0_fourpair} tests this by re-estimating the ARDL(2,4)
across all four pairings, each with and without trend, holding all other
settings identical.

\begin{table}[!ht]
	\centering
	\small
	\caption{Four-pairing fragility: ARDL(2,4) across output--capital definitions}
	\label{tab:S0_fourpair}
	\begin{tabular}{llcccc}
		\hline\hline
		Output & Capital & Trend & $\hat{\theta}$ & $\hat{\pi}_y$
		       & $1 - \sum\hat{\psi}_i$ \\
		\hline
		GVA & $K^G$ & no    & 0.750 & $-$0.089$^{*}$ & 0.111 \\
		GVA & $K^G$ & yes   & 0.565 & $+$0.013        & 0.119 \\[3pt]
		NVA & $K^G$ & no    & 0.745 & $-$0.096$^{*}$ & 0.124 \\
		NVA & $K^G$ & yes   & 0.541 & $+$0.013        & 0.132 \\[3pt]
		NVA & $K^N$ & no    & 0.807 & $-$0.052$^{*}$ & 0.063 \\
		NVA & $K^N$ & yes   & 0.129 & $-$0.029        & 0.089 \\[3pt]
		GVA & $K^N$ & no    & 0.769 & $-$0.046$^{*}$ & 0.053 \\
		GVA & $K^N$ & yes   & 0.111 & $-$0.046        & 0.077 \\
		\hline
	\end{tabular}
	\par\vspace{0.3em}
	\footnotesize
	All specifications: ARDL(2,4), Cases~II--III, dummies
	$\{d56,d74,d80\}$, $T = 65$, common deflator Py.
	$^{*}$: $\hat{\pi}_y$ significant at 10\%.
	All three dummies individually insignificant ($p > 0.5$).
\end{table}

The table splits into two regimes determined by the capital stock concept,
not the output measure. In the gross-$K$ block
($1 - \sum\hat{\psi}_i \approx 0.11$--$0.13$), trend inclusion shifts
$\hat{\theta}$ by 25--27\% and flips $\hat{\pi}_y$ from
negative-significant to positive-insignificant --- the error-correction
mechanism reverses sign. In the net-$K$ block
($1 - \sum\hat{\psi}_i \approx 0.05$--$0.09$), trend inclusion collapses
$\hat{\theta}$ by 84--86\%. Interpreted through the ARDL reparameterization
(eq.~\ref{eq:S0_lrmult}), this gradient is predicted: the denominator
$1 - \sum\hat{\psi}_i$ governs the amplification of any numerator
perturbation, and the capital stock concept determines which denominator
regime the estimation inhabits. The gross stock accumulates smoothly
(retirement rates are slow and stable), placing the divisor further from
zero; the net stock absorbs depreciation shocks that push the autoregressive
root toward the unit boundary. The fragility is a property of the ARDL's
parametric structure at this empirical site, not of the underlying
cointegrating relationship --- a distinction the VECM eigenvector estimator
will test directly.

The cointegrating evidence is fragile, surviving in one of five PSS cases
without the t-test robustness diagnostic. The trend-sensitivity is not a
pairing artifact. And Shaikh's three step dummies absorb no identifiable
structural variation across all eight specifications. These conclusions are
conditional on the maintained ARDL(2,4) specification and on the common
deflator rule. The recovered benchmark is a single point in a large
specification space --- too fragile to anchor an identification claim without
examining the space around it.


% ============================================================
% §3.4 S1: SPECIFICATION GEOMETRY
% ============================================================

\subsection{Specification Geometry}
\label{sec:S1}

% --- Opens with what S0 left unresolved ---

The S0 benchmark survives in one PSS case, with fragility governed by the
denominator regime. Whether this fragility is a property of the fixed
specification or of the method requires releasing the specification from its
fixed coordinates.

Standard practice selects a single ARDL specification by minimizing an
information criterion and treats the resulting model as the basis for
inference. This procedure embeds a researcher prior: the choice of
criterion, and the implicit weighting of fit against parsimony that each
criterion encodes, conditions the selected specification. The question is
whether the capacity anchor is recoverable as a property of a \emph{family}
of specifications occupying the Pareto frontier of fit and parametric
complexity, rather than as a property of any single winner.

The specification space is organized as a lattice:
%
\begin{equation}
	\label{eq:S1_lattice}
	\mathcal{L}_{S1} \;=\; \bigl\{\,(p,\, q,\, c,\, s)\,\bigr\},
\end{equation}
%
where $p$ is the lag order on output, $q$ the lag order on capital, $c$ the
deterministic case (Cases~I--V of \citealt{Pesaran2001}), and $s$ the
historical dummy structure. The admissible set retains only specifications
that pass the bounds test:
%
\begin{equation}
	\label{eq:S1_admissible}
	\mathcal{A}_{S1} \;=\;
	\bigl\{\, m \in \mathcal{L}_{S1} :
	F_{\text{bounds}}(m) \;\text{rejects}\; H_0 \,\bigr\}.
\end{equation}
%
Among the survivors, model comparison maps each specification into the
fit--complexity plane:
%
\begin{equation}
	\label{eq:S1_fitcomplexity}
	m \;\mapsto\; \bigl(\,k(m),\; -2\log L(m)\,\bigr),
\end{equation}
%
where information criteria appear as linear iso-penalty contours whose slope
is determined by the criterion's complexity penalty. Rather than selecting a
single tangency point, S1 identifies the fattened frontier:
%
\begin{equation}
	\label{eq:S1_frontier}
	\mathcal{F}^{(0.20)} \;=\;
	\bigl\{\, m \in \mathcal{A}_{S1} :
	\mathrm{IC}(m) \leq Q_{0.20}(\mathrm{IC}) \,\bigr\},
\end{equation}
%
the 20\% of admissible specifications with the lowest information-criterion
values --- a credible specification band rather than a single winner. Five
inferential objects are tracked across $\mathcal{F}^{(0.20)}$: the long-run
multiplier $\hat{\theta}$, the implied utilization series $\widehat{u}_t$,
the capital-stock memory share $s_K$, the deterministic case $c$, and the
dummy structure $s$. If these objects are stable across the frontier, the
capacity anchor is a property of the specification family. If they exhibit
systematic variation, the anchor is an artifact of winner selection.


% --- S1 Results ---

\subsubsection{Results}
\label{subsec:S1_results}

Table~\ref{tab:S1_lattice_stats} summarizes the specification funnel.
Of 500 candidates, 74 pass the bounds admissibility gate (14.8\%), and
$\mathcal{F}^{(0.20)}$ contains 15 specifications.
Table~\ref{tab:S1_ic_winners} reports the IC tangency points. Four of five
criteria --- AIC, BIC, HQ, and ICOMP --- converge on a single
specification: ARDL(1,3), Case~II, no dummies, with
$\hat{\theta}_{S1} = 0.565$. Only
RICOMP\footnote{Robust ICOMP: replaces the inverse Fisher information
matrix in the ICOMP penalty with the sandwich covariance estimator, making
the complexity term robust to misspecification of the error distribution
\citep{GuneyBozdoganArslan2021}. The full RICOMP pipeline is documented in
Appendix~\ref{app:ricomp}.} selects a different point: ARDL(3,3), Case~II,
no dummies, $\hat{\theta} = 0.397$. Across the frontier,
$\hat{\theta} \in [0.397,\, 0.711]$ (mean 0.573). Shaikh's ARDL(2,4)
lies outside the admissible set.

\begin{table}[!ht]
	\centering
	\small
	\caption{S1 specification lattice summary}
	\label{tab:S1_lattice_stats}
	\begin{tabular}{lrl}
		\hline\hline
		Stage & $N$ & Description \\
		\hline
		Full lattice $\mathcal{L}_{S1}$ & 500 &
			$(p,q) \in \{1,\ldots,5\}^2 \times \text{Cases I--V}
			 \times 4$ dummy sets \\
		Admissible $\mathcal{A}_{S1}$ & 74 &
			F-bounds pass at $\alpha = 0.05$ (14.8\%) \\
		Frontier $\mathcal{F}^{(0.20)}$ & 15 &
			Lowest 20\% of $-2\log L$ \\
		IC consensus (4/5) & 1 &
			ARDL(1,3), Case~II, $s_0$ \\
		RICOMP dissenter & 1 &
			ARDL(3,3), Case~II, $s_0$ \\
		Shaikh $m_0$ & --- &
			Outside admissible set \\
		\hline
	\end{tabular}
\end{table}

\begin{table}[!ht]
	\centering
	\small
	\caption{S1 information criterion winners}
	\label{tab:S1_ic_winners}
	\begin{tabular}{lcccccc}
		\hline\hline
		IC & $p$ & $q$ & Case & Dummies & $\hat{\theta}$ & $-2\log L$ \\
		\hline
		AIC   & 1 & 3 & II & none & 0.565 & $-$268.1 \\
		BIC   & 1 & 3 & II & none & 0.565 & $-$268.1 \\
		HQ    & 1 & 3 & II & none & 0.565 & $-$268.1 \\
		ICOMP & 1 & 3 & II & none & 0.565 & $-$268.1 \\
		RICOMP & 3 & 3 & II & none & 0.397 & $-$270.2 \\
		\hline
		\multicolumn{7}{l}{\footnotesize
			Shaikh $m_0$: ARDL(2,4), Case~III,
			$\{d56,d74,d80\}$, $\hat{\theta}_0 = 0.745$} \\
		\multicolumn{7}{l}{\footnotesize
			Frontier: 15 specs,
			$\hat{\theta} \in [0.397,\, 0.711]$, mean $= 0.573$} \\
	\end{tabular}
\end{table}

Three moves distinguish S1 from S0: the lag order drops from (2,4) to
(1,3), the PSS case shifts from III to II, and the dummies are rejected.
The IC consensus indicates that $\hat{\theta}_{S1} = 0.565$ is a property
of a specification family, not a contingent feature of Shaikh's lag
selection. Interpreted through the accounting framework, $\hat{\theta} < 1$
is consistent with over-mechanization: each percentage point of capital
accumulation generates less than one percentage point of productive
capacity. The RICOMP dissent --- the sandwich-corrected complexity term
penalizes collinearity near the unit-root boundary more heavily,
pulling $\hat{\theta}$ toward 0.40 --- is deferred to
Appendix~\ref{app:ricomp}.

The specification axis is confined. The conclusion is conditional on the
single-equation ARDL structure, which treats capital as weakly exogenous and
recovers the cointegrating vector through the near-singular divisor
$1 - \sum\hat{\psi}_i$. Whether the anchor survives when both conditioning
assumptions are released is the question the system estimator addresses.


% ============================================================
% §3.5 S2: SYSTEM-LEVEL VECM
% ============================================================

\subsection{System-Level VECM Replication}
\label{sec:S2}

% --- Opens with what S1 left unresolved ---

The specification axis is confined: the capacity anchor does not depend on
which information criterion selects the specification. It remains possible
that the anchor is a conditional artifact of the single-equation structure
--- that the long-run restriction would not survive if output and capital
were permitted to adjust simultaneously to the error-correction term. The
theoretical object is unchanged:
%
\begin{equation}
	\label{eq:S2_capacity}
	\ln Y_t \;-\; \theta \,\ln K_t \;=\; \ln u_t.
\end{equation}
%
What changes is the identification environment. The VECM recovers the
cointegrating vector $\beta$ as an eigenvector of a reduced-rank regression
rather than as a ratio of coefficient sums divided by a near-zero
denominator.

The specification lattice extends to six dimensions:
%
\begin{equation}
	\label{eq:S2_lattice}
	\mathcal{L}_{S2} \;=\;
	\bigl\{\,(m,\, r,\, p,\, q,\, d,\, h)\,\bigr\},
\end{equation}
%
where $m$ is the system dimension, $r$ the cointegration rank, $p$ the lag
depth, $q$ the short-run memory architecture, $d$ the deterministic branch,
and $h$ the historical shock structure. Each specification implies:
%
\begin{equation}
	\label{eq:S2_vecm}
	\Delta X_t \;=\; \alpha\beta' X_{t-1}
	\;+\; \sum_{i=1}^{p-1} \Gamma_i \,\Delta X_{t-i}
	\;+\; \text{det}(d,h)
	\;+\; \varepsilon_t,
\end{equation}
%
where $\Pi = \alpha\beta'$ with $\mathrm{rank}(\Pi) = r$, estimated via
maximum likelihood \citep{Stigler2020}. The $m = 2$ system sets
$X_t = (\ln Y_t,\, \ln K_t)'$ with $r = 1$; the $m = 3$ system extends to
$X_t = (\ln Y_t,\, \ln K_t,\, \ln e_t)'$ and admits $r = 2$, opening the
possibility of identifying reserve-army dynamics alongside the capacity
relation. Admissibility requires convergence, rank consistency with trace
and maximum-eigenvalue tests \citep{Johansen1991}, and companion-matrix
stability:
%
\begin{equation}
	\label{eq:S2_admissible}
	\mathcal{A}_{S2} \;=\;
	\bigl\{\, m \in \mathcal{L}_{S2} :
	\text{convergence} \;\wedge\;
	\text{rank} \;\wedge\;
	\text{stability}
	\,\bigr\}.
\end{equation}
%
The informational domain $\Omega_{20}$ selects the 20\% of admissible
specifications with the lowest log-likelihood loss:
%
\begin{equation}
	\label{eq:S2_omega20}
	\Omega_{20} \;=\;
	\bigl\{\, m \in \mathcal{A}_{S2} :
	{-2\log L}(m) \leq Q_{0.20}({-2\log L})
	\,\bigr\}.
\end{equation}
%
Six objects are tracked across $\Omega_{20}$: the cointegrating vector
$\beta$, the utilization series $\widehat{u}_t$, the adjustment speed
matrix $\alpha$, the capital-side memory share $s_K$, the deterministic
structure $d$, and the shock architecture $h$. If the cointegrating vector
does not rotate, the utilization series does not shift level, and the
error-correction burden does not migrate across equations as the
specification moves through $\Omega_{20}$, the capacity anchor is a system
property.

Separating the capacity transformation relation from the reserve-army
relation in the $m = 3$ system requires a rotation of the canonical
correlation space, disciplined by:
%
\begin{equation}
	\label{eq:S2_rotation}
	\hat{e}_t \;=\; \mu \;+\; \lambda\,\hat{u}_t
	\;+\; \sum_{h} \kappa_h \, DM_{h,t},
	\qquad \lambda < 0,
\end{equation}
%
where the sign restriction $\lambda < 0$ encodes the invariant that rising
capacity utilization exerts a negative cyclical effect on exploitation
pressure.


% --- S2 Results ---

\subsubsection{Results}
\label{subsec:S2_results}

Table~\ref{tab:S2_ic_winners} reports the IC-selected VECM specifications.
In the bivariate system, all four criteria select lag depth $p = 1$ with
$\hat{\theta}_{S2} \in [0.870,\, 0.900]$; the $\Omega_{20}$ mean is
$\hat{\theta} = 0.803$. Table~\ref{tab:S2_lattice_stats} summarizes the
funnel: of 48 candidates, 21 pass the triple gate (43.8\%).

The trivariate system produces a sign split: AIC and HQ recover
$\hat{\theta} = -1.20$ under the unrestricted-constant specification
($d_3$), while BIC and ICOMP recover $\hat{\theta} \approx 0.93$ under
$d_0$. The sign degeneracy does not disqualify the trivariate system --- it
signals that the reduced-form canonical correlations have not been rotated
into economically identified vectors. Whether the rotation
(eq.~\ref{eq:S2_rotation}) resolves the degeneracy is tested alongside the
regime-break analysis, where the trivariate system re-enters under
structural identification constraints.

\begin{table}[!ht]
	\centering
	\small
	\caption{S2 IC winners: bivariate and trivariate VECM}
	\label{tab:S2_ic_winners}
	\begin{tabular}{clccccc}
		\hline\hline
		$m$ & IC & $p$ & Det.\ & Dummies & $\hat{\theta}$ & $-2\log L$ \\
		\hline
		2 & AIC   & 1 & $d_2$ & $h_1$ & 0.885 & $-$572.5 \\
		2 & BIC   & 1 & $d_0$ & $h_0$ & 0.870 & $-$560.4 \\
		2 & HQ    & 1 & $d_2$ & $h_1$ & 0.885 & $-$572.5 \\
		2 & ICOMP & 1 & $d_0$ & $h_1$ & 0.900 & $-$566.5 \\
		\hline
		3 & AIC   & 3 & $d_3$ & $h_0$ & $-$1.195 & $-$802.3 \\
		3 & BIC   & 3 & $d_0$ & $h_0$ & 0.930     & $-$783.9 \\
		3 & HQ    & 3 & $d_3$ & $h_0$ & $-$1.195 & $-$802.3 \\
		3 & ICOMP & 3 & $d_0$ & $h_0$ & 0.930     & $-$783.9 \\
		\hline
	\end{tabular}
	\par\vspace{0.3em}
	\footnotesize
	Bivariate $\Omega_{20}$ mean: $\hat{\theta} = 0.803$.
	Trivariate sign split: reduced-form degeneracy, rotation pending.
\end{table}

\begin{table}[!ht]
	\centering
	\small
	\caption{S2 VECM lattice summary}
	\label{tab:S2_lattice_stats}
	\begin{tabular}{clrl}
		\hline\hline
		$m$ & Stage & $N$ & Description \\
		\hline
		2 & Full lattice & 48 & $(p,d,h)$ grid \\
		2 & Admissible & 21 & Johansen rank $\geq 1$ (43.8\%) \\
		2 & $\Omega_{20}$ & 5 & Lowest 20\% of $-2\log L$ \\
		2 & IC band & $[0.87,\, 0.90]$ & All 4 ICs \\
		\hline
		3 & Admissible & low & $\sim$2\% admissibility \\
		3 & Sign split & 2/4 &
			AIC/HQ: $\hat{\theta} < 0$;
			BIC/ICOMP: $\hat{\theta} > 0$ \\
		\hline
	\end{tabular}
\end{table}

Both estimators indicate $\hat{\theta} < 1$: over-mechanization is robust
across the single-equation and system frameworks. The point estimates do not
coincide: $\hat{\theta}_{S1} \approx 0.57$ versus
$\hat{\theta}_{S2} \approx 0.88$. Interpreted through the measurement
framework, this gap is the empirical counterpart of the denominator
fragility diagnosed in the four-pairing exercise
(Table~\ref{tab:S0_fourpair}): the ARDL recovers $\hat{\theta}$ through
$\sum\hat{\varphi}_j / (1 - \sum\hat{\psi}_i)$, where the divisor is
approximately 0.12; the VECM eigenvector estimator bypasses that divisor.
The level discrepancy is consistent with method-induced depression of the
ARDL point estimate, not with a structural difference in the cointegrating
relationship.

The trivariate sign split indicates that rank~1 is the identification
boundary of the unrestricted bivariate system. This conclusion is
conditional on treating the unrestricted VECM as the identification
environment; it is possible that the degeneracy reflects entanglement of
the capacity and reserve-army relations that the sign restriction
$\lambda < 0$ can resolve. Two questions pass forward: whether
$\hat{\theta}$ is stable across the Fordist/Post-Fordist periodization,
and whether the trivariate degeneracy resolves under structural rotation.


% ============================================================
% §3.6 CROSS-STAGE SUMMARY
% ============================================================

\subsection{Cross-Stage Summary}
\label{subsec:cross_stage}

Table~\ref{tab:cross_theta} collects the $\hat{\theta}$ trajectory. The
qualitative invariant --- $\hat{\theta} < 1$, consistent with
over-mechanization --- holds across every admissible specification in every
stage. The quantitative variation is governed by two identified sources: the
denominator regime (gross vs.\ net capital stock, diagnosed in S0) and the
estimator class (ARDL near-singular divisor vs.\ VECM eigenvector, diagnosed
in S2). Neither source indicates a different economic relationship; both are
properties of the parametric recovery of the long-run multiplier at this
empirical site.

\begin{table}[!ht]
	\centering
	\small
	\caption{Cross-stage $\hat{\theta}$ trajectory}
	\label{tab:cross_theta}
	\begin{tabular}{llcl}
		\hline\hline
		Stage & Estimator & $\hat{\theta}$ & Note \\
		\hline
		S0 & ARDL(2,4), Case~II & 0.745 &
			Sole passing case \\
		S0 & ARDL(2,4), Cases~IV--V & 0.541 &
			Trend collapse ($-$27\%) \\
		S0 & Four pairings, gross $K$ & $[0.54,\, 0.75]$ &
			GVA/NVA $<$ 1 pp \\
		S0 & Four pairings, net $K$ & $[0.11,\, 0.81]$ &
			Denominator-amplified \\
		\hline
		S1 & ARDL(1,3), Case~II & 0.565 &
			4/5 ICs; no dummies \\
		S1 & RICOMP dissenter & 0.397 &
			Sandwich penalty \\
		S1 & $\mathcal{F}^{(0.20)}$ frontier & $[0.40,\, 0.71]$ &
			15 specs \\
		\hline
		S2 & VECM $m\!=\!2$, IC band & $[0.87,\, 0.90]$ &
			Eigenvector \\
		S2 & VECM $m\!=\!2$, $\Omega_{20}$ & 0.803 &
			Domain mean \\
		S2 & VECM $m\!=\!3$ & sign split &
			Rotation pending \\
		\hline
	\end{tabular}
	\par\vspace{0.3em}
	\footnotesize
	All stages: $\hat{\theta} < 1$ (over-mechanization). \\
	S1$\to$S2 gap: denominator fragility (ARDL divisor vs.\ VECM eigenvector). \\
	$m=3$ sign degeneracy: structural rotation deferred.
\end{table}

\noindent
The confined bivariate specification ($m = 2$, $p = 1$, Case~II) and the
unresolved trivariate degeneracy pass forward to the regime-break analysis,
which tests whether $\hat{\theta}$ is stable across the
Fordist/Post-Fordist periodization and whether the sign restriction
$\lambda < 0$ resolves the trivariate entanglement under structural
identification constraints.


\section{References}
\bibliographystyle{apalike}
\bibliography{references}
\label{references}

%%%%%%%%%%%%%%%%%%%%%%%%%%%%%%%%%%
%%%%  Appendix %%%%%%%%%%%%%%%%%%%
%%%%%%%%%%%%%%%%%%%%%%%%%%%%%%%%%%
\pagebreak
\appendix


\section{Data Glossary and Construction Ledger}
\label{app:data_ledger}

% ---------------------------------------------------------------
% A.1 — Series glossary
% ---------------------------------------------------------------

\subsection{Series Glossary}

\begin{table}[!ht]
	\centering
	\small
	\setlength{\tabcolsep}{3pt}
	\renewcommand{\arraystretch}{1.15}
	\caption{Replication series: symbols, sources, and transformations}
	\label{tab:app_glossary}
	\begin{tabular}{@{}>{\raggedright\arraybackslash}p{0.20\textwidth}
			>{\raggedright\arraybackslash}p{0.26\textwidth}
			>{\raggedright\arraybackslash}p{0.28\textwidth}
			>{\raggedright\arraybackslash}p{0.18\textwidth}@{}}
		\hline
		Symbol / Code & Variable & Source (table, line) & Transformation \\
		\hline
		$\text{GVAcorp}_t$ 
		& Corporate gross value added (corrected)
		& NIPA 1.14 (L1) + imputed-interest adj.\ (Table~7.11)
		& nominal \\[3pt]
		$\text{KGCcorp}_t$ 
		& Corporate gross capital stock (GPIM-adjusted)
		& Fixed Asset 6.1 (L9); IRS Statistics of Income; GPIM recursion
		& nominal \\[3pt]
		$\text{KNCcorp}_t$ 
		& Corporate net capital stock (GPIM-adjusted)
		& same as above
		& nominal \\[3pt]
		$\text{KNCcorpbea}_t$
		& Corporate net capital stock (BEA 2011 official)
		& Fixed Asset 6.1 (L9)
		& nominal \\[3pt]
		$p_t$ 
		& Common deflator
		& NIPA 1.1.9 (L1)
		& index \\[3pt]
		$y_t$
		& Real log output
		& constructed
		& $\ln(\text{GVAcorp}_t/p_t)$ \\[3pt]
		$k_t$
		& Real log capital
		& constructed
		& $\ln(\text{KGCcorp}_t/p_t)$ \\[3pt]
		$\text{Profshcorp}_t$ 
		& Corporate profit share (corrected)
		& NIPA 1.14 (L8, L1$-$L2) + imputed-interest adj.
		& ratio \\
		\hline
	\end{tabular}
	\par\vspace{0.4em}
	\footnotesize\textit{Note:} All BEA and NIPA data downloaded in the
	replication vintage. Shaikh's original series used data downloaded in
	2011 \citep[Appendix~6.7, fn.~1]{Shaikh2016}; discrepancies attributable
	to subsequent BEA revisions are possible and noted where material.
	\par\vspace{0.2em}
	\footnotesize\textit{Source:}
	\href{https://www.realecon.org/data}{Real Economic Analysis Data Tables}
\end{table}

% ---------------------------------------------------------------
% A.2 — Imputed-interest correction
% ---------------------------------------------------------------

\subsection{Imputed-Interest Correction}

The NIPA corporate value added figure includes imputed financial
intermediation services as a cost to non-financial corporations. The
corporate imputed-interest adjustment removes this component:
%
\begin{equation}
	\label{eq:app_impint}
	\text{CorpImpIntAdj}_t
	\;\equiv\;
	-\,\text{BankMonIntPaid}_t
	\;-\; \text{CorpNFNetImpIntPaid}_t,
\end{equation}
%
where $\text{BankMonIntPaid}_t$ is NIPA Table~7.11
(lines 4+44+73 minus lines 28+52+91) and
$\text{CorpNFNetImpIntPaid}_t$ is NIPA Table~7.11 (line~74 minus line~53).
In 2009 the adjustment raises corporate net operating surplus by
approximately 10\,\%.

% ---------------------------------------------------------------
% A.3 — Corporate output: GVAcorp
% ---------------------------------------------------------------

\subsection{Corrected Corporate Output: \texorpdfstring{$\text{GVAcorp}_t$}{GVAcorp}}

\begin{equation}
	\label{eq:app_gva}
	\text{GVAcorp}_t
	\;\equiv\;
	\text{GVAcorpnipa}_t \;+\; \text{CorpImpIntAdj}_t,
\end{equation}
%
where $\text{GVAcorpnipa}_t$ is NIPA Table~1.14 (line~1) and
$\text{CorpImpIntAdj}_t$ is defined in equation~\eqref{eq:app_impint}.

% ---------------------------------------------------------------
% A.4 — Corporate capital stock: KGCcorp (GPIM)
% ---------------------------------------------------------------

\subsection{Reconstructed Corporate Capital Stock: \texorpdfstring{$\text{KGCcorp}_t$}{KGCcorp}}

\paragraph{GPIM accumulation rule.}
The current-cost capital stock is accumulated as:
%
\begin{equation}
	\label{eq:app_gpim}
	K_t \;=\; IG_t \;+\; z_t^* \cdot K_{t-1},
	\qquad
	z_t^* \;\equiv\; (1 - z_t)\,\frac{p^K_t}{p^K_{t-1}},
\end{equation}
%
where $IG_t$ is nominal gross investment, $z_t$ is the period depreciation
rate, and $z_t^*$ is the survival-and-revaluation factor that reflates the
surviving stock to current replacement cost.

\paragraph{Three adjustments (applied before the recursion).}

\begin{enumerate}
	\item \textbf{Depletion rates.} BEA 1993 finite service lives replace
	the BEA 2011 infinite-lives assumption, implying higher depreciation
	and lower retirement rates consistent with physical exhaustion.
	
	\item \textbf{Initial values.} BEA 1993 initial values anchored to
	1925 replace the BEA 2011 starting values. The 1993 values are lower,
	producing faster accumulation growth and convergence to the official
	path by approximately 1977.
	
	\item \textbf{Great Depression correction.} The IRS corporate book
	value index \citep[Series~V~115, pp.~924--926]{Census1975} is applied
	over 1925--1947, rescaling the BEA historical stock by the IRS/BEA
	ratio each year. The GPIM recursion resumes from 1947. Net effect:
	$\text{KNCcorp}_t$ is approximately 28\,\% lower in 1947 than the BEA
	2011 official measure.
\end{enumerate}

\paragraph{Inventories.}
IRS corporate balance sheet inventory ratios are extrapolated over
1947--1989 via a linear fit to the 1990--2011 IRS/BEA ratio, rescaled to
the adjusted GPIM stocks, and added to $\text{KGCcorp}_t$. Source:
IRS Statistics of Income 1926--2011; FRB Flow of Funds FL105015205.A for
non-financial corporate inventories in current cost.

\paragraph{Accuracy.}
The average ratio of the GPIM approximation to the BEA current-cost stock
over 1947--2005 is 99.60\,\%.

\noindent Full recursion inputs and depletion-rate construction are in
\citet[Appendix~6.7, Section~V; Appendix Table~6.8.II.3]{Shaikh2016}.

% ---------------------------------------------------------------
% A.5 — Common deflator
% ---------------------------------------------------------------

\subsection{Common Deflator: \texorpdfstring{$p_t$}{p\_t}}

The replication uses the GDP implicit price deflator (BEA NIPA
Table~1.1.9, line~1) as $p_t$. This is a workaround: Shaikh's
Appendix~6.8 tables do not tabulate the output-side deflator used in the
corporate utilization exercise (see \S\ref{sec:ch3_data_gap}). The
transformed series are:
%
\begin{equation}
	\label{eq:app_deflator}
	y_t \;\equiv\; \ln\!\left(\frac{\text{GVAcorp}_t}{p_t}\right),
	\qquad
	k_t \;\equiv\; \ln\!\left(\frac{\text{KGCcorp}_t}{p_t}\right).
\end{equation}

% ---------------------------------------------------------------
% A.6 — Corrected corporate profit share: Profshcorp
% ---------------------------------------------------------------

\subsection{Corrected Corporate Profit Share: \texorpdfstring{$\text{Profshcorp}_t$}{Profshcorp}}

The same imputed-interest adjustment applied to $\text{GVAcorp}_t$ is
applied to both the numerator and denominator of the profit share:
%
\begin{align}
	\text{NOScorp}_t
	&\;\equiv\; \text{NOScorpnipa}_t \;+\; \text{CorpImpIntAdj}_t,
	\label{eq:app_nos} \\
	\text{VAcorp}_t
	&\;\equiv\; \text{VAcorpnipa}_t \;+\; \text{CorpImpIntAdj}_t,
	\label{eq:app_va} \\
	\text{Profshcorp}_t
	&\;\equiv\; \frac{\text{NOScorp}_t}{\text{VAcorp}_t},
	\label{eq:app_profshcorp}
\end{align}
%
where $\text{NOScorpnipa}_t$ is NIPA Table~1.14 (line~8) and
$\text{VAcorpnipa}_t$ is NIPA Table~1.14 (line~1 minus line~2).
$\text{Profshcorp}_t$ is constructed in the S0 data environment but does
not enter the S0 or S1 identification objects; its first use is as the
empirical proxy for $\ln e_t$ in the S2 $m=3$ VECM system.

\section{Electric Motor Utilization: Physical Measures Lineage}
\label{app:emu_lineage}

The Electric Motor Utilization index represents the founding device of the
physical-intensity tradition in capacity measurement. Unlike survey-based
approaches, which elicit managerial reports of operating rates relative to a
self-defined normal, the EMU methodology grounds the utilization ratio in an
engineering relationship between installed motor capacity and observed
electricity consumption. Table~\ref{tab:EMU_proxy_CU} documents the
methodology, data sources, and successive extensions of this lineage from
\citet{Foss1963} through \citet{Shaikh1992,Shaikh2016} to
\citet{Nikiforos2016}, who formalized the physical-intensity critique as a
micro-founded theoretical result and proposed the Average Workweek of Capital
as the contemporary analogue.

\begin{table}[!h]
	\centering
	\caption{Electric Motor Utilization (EMU) as a Proxy for Capacity
		Utilization: Physical Measures Lineage}
	\label{tab:EMU_proxy_CU}
	\small
	\begin{tabular}{p{3.2cm} p{10.5cm}}
		\toprule
		Feature & Description \\
		\midrule
		
		Time Span &
		1929--1954 (manufacturing); extended to 1899--1963 by
		\citet{Shaikh1992}. \\[4pt]
		
		Data \& Methodology &
		Calculates potential motor output (in kWh) from horsepower of
		installed electric motors $\times$ 8{,}760 hours/year $\times$
		0.746 kWh per hp-hour $\div$ 0.9 (efficiency adjustment). Actual
		electricity consumption is then compared to this theoretical
		maximum. Primary sources: \citet{USfpc1945};
		\citet{UScensus1966,UScensus1967,UScensus1975};
		\citet{USdc1986}. \\[4pt]
		
		Main Results &
		Substantial long-term increases in motor utilization from 1929 to
		1954: utilization rates rose from 15.9\% (1929) to 20.9\% (1954),
		indicating rising intensity in capital use even as labour hours
		declined \citep{Foss1963}. \\[4pt]
		
		Critique of FRB &
		Indirect. Suggests that survey-based measures overlook mechanical
		intensity and changes in input use; implicitly argues for a more
		objective, physically grounded metric capable of registering
		greater long-run variation than survey-elicited operating rates. \\[4pt]
		
		Extensions by \citet{Shaikh1992} &
		Adapts and extends the EMU method using Census data and
		\citet{Christensen1969} to build a continuous series from
		1899--1963, adjusting pre-1939 estimates and interpolating from
		industry-level motor data. Spliced with McGraw-Hill survey data to
		produce a utilization measure covering 1947--1985. The extended
		series diverges from FRB in amplitude, trend, and cyclical
		pattern, capturing episodes the FRB series suppresses
		\citep[Appendix~6.6]{Shaikh2016}. \\[4pt]
		
		Formalization by \citet{Nikiforos2016} &
		Re-introduces the Average Workweek of Capital (AWW) --- hours per
		week that capital equipment is operated --- as the contemporary
		physical-intensity analogue to the EMU ratio. Formalizes the
		critique of survey-based measures through the plant-manager
		thought experiment: because adding or dropping shifts changes the
		normal against which managers report, the survey denominator moves
		with the numerator. Stationarity of the FRB series is thereby a
		construction artefact, not a property of the underlying utilization
		dynamics. AWW-based estimates exhibit non-stationary behaviour
		consistent with the EMU lineage. \\[4pt]
		
		Primary Sources &
		\citet{USfpc1945}; \citet{UScensus1966}; \citet{UScensus1967};
		\citet{UScensus1975}; \citet{USdc1986}; \citet{Christensen1969};
		\citet{Foss1963}; \citet{Shaikh1992,Shaikh2016};
		\citet{Nikiforos2016}. \\
		
		\bottomrule
	\end{tabular}
	\smallskip
	
	\noindent\textit{Source: Own elaboration based on \citet{Foss1963},
		\citet{Shaikh1992,Shaikh2016}, and \citet{Nikiforos2016}.}
\end{table}

\end{document}